\documentclass[11pt]{article}

\usepackage[margin=1in]{geometry}
\usepackage{amsmath,amssymb}
\usepackage{booktabs}
\usepackage{hyperref}
\usepackage[numbers,sort&compress]{natbib}

\title{CASPO: Learned Prefix Value Modeling for\\
       Stepwise Credit Assignment in Reasoning RL}
\author{Anonymous}
\date{}

\newcommand{\caspo}{\textsc{CASPO}}
\newcommand{\ppo}{PPO}
\newcommand{\grpo}{GRPO}
\newcommand{\ipvrm}{IPVRM}
\newcommand{\vineppo}{VinePPO}

\begin{document}
\maketitle

\begin{abstract}
We propose \caspo{} (Credit Assignment Step Policy Optimization), a method for
reinforcement learning of mathematical reasoning that obtains step-level
credit signals from a learned prefix-value model rather than from on-policy
Monte Carlo prefix rollouts. \caspo{} retains the \vineppo{} step-segmentation
and step-TD advantage structure but replaces \vineppo{}'s $K{=}9$ rollouts at
each non-terminal prefix with a single forward pass through an
\ipvrm{}-style cumulative-log-ratio prefix value model, optionally fine-tuned
online during policy training under an ADB+DLW correction. The result is a
strict drop-in for \vineppo{} that uses the same PPO clipped objective, the
same KL penalty, the same response segmentation, and the same global PPO
minibatch, but eliminates the per-step Monte Carlo expansion that dominates
\vineppo{}'s wall-clock cost. We compare four methods (\ppo{}, \grpo{},
\vineppo{}, \caspo{}) under one trainer, one verifier, one rollout engine, and
one set of hyperparameters, and report wall-clock and accuracy on the MATH
benchmark for both Rho-1B and DeepSeekMath-7B. Throughout we keep the
algorithmic contract identical except for the source of the per-step value
$V(s_t)$, isolating the credit-signal change from confounders. Empirical
results are reported in Section~\ref{sec:results} after the corresponding
training and evaluation runs complete.
\end{abstract}

\section{Introduction}
\label{sec:intro}

Reinforcement learning has emerged as the dominant approach for improving the
mathematical reasoning of large language models, but the dominant signal is
sparse: a verifier returns 0/1 only at the end of a long chain of thought,
leaving the optimizer to assign credit across hundreds of intermediate tokens.
\ppo{} \citep{schulman2017ppo} addresses this by broadcasting a sequence-level
advantage uniformly across response tokens; \grpo{} \citep{deepseekmath2024}
refines the variance term using a group-relative baseline. Neither method
distinguishes a correct intermediate inference from a lucky guess that
nonetheless reached the right final answer.

\vineppo{} \citep{vineppo2024} addresses this gap by introducing
\emph{step-level} value estimates: the response is segmented into reasoning
steps, and the value at each step boundary $s_t$ is estimated by sampling $K$
on-policy continuations from $s_t$ and averaging their terminal verifier
rewards. This produces well-calibrated step values that bear out
empirically---\vineppo{} reports substantial improvements over \ppo{} and
\grpo{} on MATH and GSM8K---but at a cost: each non-terminal step boundary
requires $K$ additional generations, multiplying wall-clock time by roughly
the average step count per response.

\caspo{} replaces those Monte Carlo prefix rollouts with a single forward
pass through a \emph{learned} prefix value model
$V_\phi(s_t)$. Following \ipvrm{} \citep{ipvrm2026}, $V_\phi$ is parameterized
as a cumulative log-ratio between a trainable prefix policy $\pi_\phi$ and a
frozen reference $\pi_{\mathrm{ref}}$, and trained from trajectory-level
correctness labels using a margin BCE objective. \caspo{} keeps every other
\vineppo{} component fixed: the same LaTeX-aware step splitter, the same
step-TD advantage construction, the same clipped PPO objective, and the same
KL penalty. Step-level credit is therefore extracted from $V_\phi$ in
$\mathcal{O}(1)$ forward passes per response rather than $\mathcal{O}(K \cdot
\text{steps})$ generations, and online updates of $V_\phi$ during policy
training mitigate distribution shift.

\paragraph{Contributions.} We contribute:
(i) \caspo{}: a method that drops \vineppo{}'s prefix Monte Carlo step in
favor of an \ipvrm{}-style learned prefix value model, including the step-TD
construction and an online update path with ADB and DLW corrections;
(ii) a family of three principled step-TD advantage formulations
(\emph{value}, \emph{prob}, \emph{logprob}) derived from the Bradley--Terry
interpretation of $V_\phi$ under the IPVRM BCE loss, each operating on a
different scale of the prefix value (raw cumulative log-ratio, implied
correctness probability, or its log) and motivated by a different
bias--variance regime;
(iii) a frozen reward-model variant (\caspo{}-frozen-RM) that disables online
value updates and isolates the cost/benefit of online value drift correction;
(iv) a unified comparison stack that implements \ppo{}, \grpo{}, \vineppo{},
and \caspo{} under one trainer, one verifier, and one rollout engine,
controlling everything except the algorithmic credit-signal choice; and
(v) implementation infrastructure for full-finetune 7B-scale RL with
FSDP-sharded trainer and rank-local vLLM generation, including a CUDA-IPC
weight-sync path that is collective with \texttt{summon\_full\_params}.

\paragraph{Implementation provenance.} The codebase is an adaptation and
extension of the public \vineppo{} setup \citep{vineppo2024,vineppocode}. The
Rho-1B and DeepSeekMath-7B configurations, prompt templates, rollout shapes,
PPO hyperparameters, evaluation protocol, and LaTeX-aware step segmentation
are matched to or derived from the upstream paper and reference repository so
that any wall-clock or accuracy gap between methods reflects the algorithmic
change rather than infrastructure.

\section{Related Work}
\label{sec:related}

\paragraph{Policy gradient with sparse rewards.}
\ppo{} \citep{schulman2017ppo} optimizes a clipped policy-gradient surrogate
to bound the per-step policy update. Applied to language models with
verifier-only terminal rewards, the standard practice is to broadcast the
sequence-level advantage uniformly across response tokens and add a
KL-to-reference penalty to prevent the policy from drifting far from the SFT
initialization.

\paragraph{Group-relative baselines.}
\grpo{} \citep{deepseekmath2024} samples $G$ responses per prompt and uses
the within-group reward statistics as the advantage baseline. This avoids
training a separate value network entirely, at the cost of treating every
intermediate token in a correct response as equally important.

\paragraph{Step-level credit assignment via Monte Carlo prefix values.}
\vineppo{} \citep{vineppo2024} introduces stepwise value estimates by
sampling $K$ on-policy continuations from each non-terminal prefix and
averaging the verifier outcomes of those continuations. Step advantages are
then formed as temporal-difference differences between consecutive prefix
values. The reasoning-credit gain is well demonstrated on MATH, but the
method's rollout multiplier scales as $K$ times the average number of steps
per response.

\paragraph{Learned process supervision.}
A line of work trains a separate process reward model on human-annotated or
self-supervised step-correctness labels. \caspo{}'s prefix value model is in
this family but constructed differently: rather than scoring step text
directly, it accumulates token-level $\log
[\pi_\phi/\pi_{\mathrm{ref}}]$ contributions, following \ipvrm{}
\citep{ipvrm2026}. This aligns prefix scoring with the same density that
generated the response and avoids requiring a separate textual annotation.

\paragraph{Implicit prefix value learning (\ipvrm{}).}
\ipvrm{} \citep{ipvrm2026} formalizes prefix values as cumulative log-ratios
of a trainable prefix policy versus a frozen reference policy, trained from
trajectory-level outcome labels using a margin BCE. \ipvrm{}'s reported
setting uses LoRA adapters on a 7B base; we adopt the same parameterization
but perform full-model updates, with corresponding adjustments to the online
learning rate.

\paragraph{Generation backends and weight sync.}
We use vLLM \citep{kwon2023vllm} for high-throughput rollout and prefix
caching. The implementation uses CUDA-IPC weight transfer between trainer
ranks and rank-local vLLM engines; for the FSDP-sharded 7B trainer, the IPC
sync is implemented as a collective \texttt{summon\_full\_params} block over
the FSDP-wrapped policy followed by per-tensor IPC handle exchange with each
rank's vLLM \texttt{EngineCore}.

\section{Background}
\label{sec:bg}

\subsection{Reasoning RL setup}
A trajectory consists of a problem prompt $x$ and a model response
$y = (y_1, \dots, y_T)$ generated autoregressively. A verifier
$R(x,y) \in \{0,1\}$ returns 1 if and only if the boxed final answer in $y$
matches the ground truth. The training objective is to maximize
$\mathbb{E}_{x \sim \mathcal{D}, y \sim \pi_\theta(\cdot \mid x)} R(x,y)$
subject to a KL constraint to a frozen reference policy $\pi_{\mathrm{ref}}$,
which we instantiate as the SFT initialization.

\subsection{Step segmentation}
We split each response into reasoning steps using the LaTeX-aware splitter of
\vineppo{} \citep{vineppo2024,vineppocode}, which identifies natural
mathematical step boundaries (display equations, aligned blocks, sentence
terminators in prose). Let $0 = t_0 < t_1 < \dots < t_S = T$ be the resulting
boundary token indices and let $s_t \triangleq y_{<t}$ denote the prefix
state at boundary $t$.

\subsection{Step-TD advantage}
Given a per-step value function $V(s_t)$, both \vineppo{} and \caspo{}
construct the step-TD advantage
\begin{equation}
\label{eq:adv}
A(s_t) = r_t + \gamma\, V(s_{t+1}) - V(s_t),
\qquad
r_t = \begin{cases} R(x,y) & t = S, \\ 0 & t < S, \end{cases}
\end{equation}
with $\gamma = 1$. \caspo{} obtains $V$ from a learned model
$V_\phi$; \vineppo{} obtains $V$ by averaging $K$ Monte Carlo continuations.
The downstream PPO objective and KL penalty are identical for both methods;
only the source of $V$ differs.

\subsection{PPO objective}
We use the clipped PPO objective
\begin{equation}
\label{eq:ppo}
\mathcal{L}_{\mathrm{clip}}(\theta) =
-\,\mathbb{E}_t\!\left[
\min\!\big(
\rho_t(\theta)\, A_t,\;
\mathrm{clip}\!\big(\rho_t(\theta), 1-\epsilon, 1+\epsilon\big)\, A_t
\big)
\right],
\end{equation}
with $\rho_t(\theta) = \pi_\theta(y_t \mid x, y_{<t}) / \pi_{\theta_{\mathrm{old}}}(y_t \mid x, y_{<t})$
and $\epsilon = 0.2$. A KL-to-reference term is added with the $k_3$
estimator:
\begin{equation}
\label{eq:kl}
\mathrm{KL}_{k_3}(\pi_\theta \,\Vert\, \pi_{\mathrm{ref}}) \approx
\mathbb{E}_t\!\left[
\exp(\delta_t) - \delta_t - 1
\right], \quad
\delta_t = \log \pi_{\mathrm{ref}}(y_t) - \log \pi_\theta(y_t).
\end{equation}
The advantages $A_t$ in \eqref{eq:ppo} are (a) terminal-broadcast for
\ppo{}/\grpo{} and (b) the step-TD advantages of \eqref{eq:adv} broadcast over
the tokens in each step span for \vineppo{}/\caspo{}.

\section{Method: \caspo{}}
\label{sec:method}

\subsection{Prefix value model}
Following \ipvrm{} \citep{ipvrm2026}, we define
\begin{equation}
\label{eq:vphi}
V_\phi(s_t) \;=\;
\beta \sum_{i < t}
\log \frac{\pi_\phi(y_i \mid x, y_{<i})}{\pi_{\mathrm{ref}}(y_i \mid x, y_{<i})},
\end{equation}
where $\pi_\phi$ is a trainable prefix policy initialized from the same SFT
checkpoint family as the base policy and $\pi_{\mathrm{ref}}$ is the same
frozen reference used in the KL term in \eqref{eq:kl}. The scale $\beta$ is a
fixed scalar; we use $\beta = 10$ throughout.

\subsection{Phase-1: offline value training}
We train $V_\phi$ on a fixed corpus of (prompt, response, outcome) triples
collected by sampling $G$ responses per prompt from the SFT policy at
generation temperature $1.0$, scoring with the verifier, and discarding
prompts whose $G$ outcomes are all 0 or all 1 (the mixed-outcome filter,
following \ipvrm{}). The training loss is a margin-BCE between paired
correct/incorrect prefixes:
\begin{equation}
\label{eq:ipvrmloss}
\mathcal{L}_{\mathrm{IPVRM}}(\phi) =
\mathbb{E}_{(x, y^+, y^-)}\!\left[
\log\!\big(1 + \exp(-(V_\phi(y^+) - V_\phi(y^-) - m))\big)
\right],
\end{equation}
with margin $m = 5$. We train for three epochs with full-model updates and a
small learning rate ($5 \times 10^{-7}$).

\subsection{Phase-2: step-TD policy optimization}
At each PPO outer step, the trainer samples $G$ responses per prompt for
$P$ prompts, segments each response, and constructs step values
$\{V_\phi(s_t)\}$ in a single forward pass through $V_\phi$ using
\eqref{eq:vphi}. Step advantages are then assembled from \eqref{eq:adv} and
broadcast over the tokens within each step span before applying the PPO
objective in \eqref{eq:ppo} with KL penalty in \eqref{eq:kl}. The standard
\caspo{} run optionally co-trains $V_\phi$ during policy optimization (next
section).

\subsection{Online value updates with ADB+DLW}
Distribution drift between the SFT-init prefix policy and the actively
optimized policy can degrade $V_\phi$. We therefore allow online updates of
$V_\phi$ from the same rollout batch used for PPO. The on-policy update uses
the same margin-BCE loss \eqref{eq:ipvrmloss} but with two corrections from
\ipvrm{}:
\begin{itemize}
\item \emph{ADB} (advantage debias): the BCE boundary is shifted per-prompt
by $\log\!\big[\mu/(1-\mu)\big]$, where $\mu$ is the empirical positive
fraction over the prompt's $G$ rollouts, so the value model learns
\emph{deviations} from the prompt's SFT difficulty rather than its absolute
difficulty.
\item \emph{DLW} (difficulty-level weighting): the BCE term is weighted by
$\mu (1-\mu)$, downweighting prompts where outcomes are nearly all-correct
or all-incorrect.
\end{itemize}
The online value learning rate is set to $10^{-6}$, an order of magnitude
below the offline rate to control drift.

\subsection{Three step-TD advantage formulations from $V_\phi$}
\label{sec:transforms}

A central modeling choice in \caspo{} is \emph{which} statistic of $V_\phi$
enters the step-TD difference. We derive and propose three principled
formulations from how the IPVRM loss \eqref{eq:ipvrmloss} parameterizes
correctness. Recall that the BCE margin loss treats $V_\phi$ as a
Bradley--Terry score: the implied per-prefix correctness probability is
\begin{equation}
\label{eq:ptilde}
\widetilde{p}(s_t) \;\triangleq\; \sigma(V_\phi(s_t)) \;\approx\; \mathbb{P}\!\left[R(x,y)=1 \mid s_t\right],
\end{equation}
where $\sigma$ is the logistic function. Equation~\eqref{eq:ptilde} is the
natural way to convert an unbounded cumulative-log-ratio score into a
calibrated probability, and it is the quantity the BCE objective directly
optimizes.

Equation~\eqref{eq:adv} is then well-defined on three distinct scales of
$V_\phi$, each with a different statistical interpretation. We propose all
three as named formulations of \caspo{}:

\begin{description}
\item[\emph{value}] (raw cumulative log-ratio):
\begin{equation}
\label{eq:advV}
A^{V}(s_t) \;=\; r_t \,+\, \gamma\, V_\phi(s_{t+1}) \,-\, V_\phi(s_t).
\end{equation}
\item[\emph{prob}] (Bradley--Terry probability):
\begin{equation}
\label{eq:advP}
A^{P}(s_t) \;=\; r_t \,+\, \gamma\, \widetilde{p}(s_{t+1}) \,-\, \widetilde{p}(s_t)
            \;=\; r_t \,+\, \gamma\, \sigma(V_\phi(s_{t+1})) \,-\, \sigma(V_\phi(s_t)).
\end{equation}
\item[\emph{logprob}] (log Bradley--Terry probability):
\begin{equation}
\label{eq:advLP}
A^{LP}(s_t) \;=\; r_t \,+\, \gamma\, \log \widetilde{p}(s_{t+1}) \,-\, \log \widetilde{p}(s_t).
\end{equation}
\end{description}

In all three variants, the terminal verifier reward $r_T = R(x,y)$ enters
identically; the only difference is the function applied to $V_\phi$ before
the TD difference.

\paragraph{Intuition.}
The three transforms have different scales and different bias--variance
profiles:

\begin{itemize}
\item \emph{value}: The TD operates on the unbounded log-ratio scale. A
single token whose log-ratio swings by, say, $-10$ produces an unbounded
$\Delta V_\phi$, which carries a strong gradient signal but also amplifies
miscalibrated regions of the prefix value model (e.g.\ tokens far from the
ref-policy distribution). High signal-to-noise where $V_\phi$ is
well-calibrated, but heavier tails when it is not. Standardization +
clipping (Section~\ref{sec:method}, advantage clip $\pm 3$) bounds the
magnitude empirically, but does not reshape the underlying scale.

\item \emph{prob}: The TD operates on the probability scale, so
$|A^P(s_t)| \leq 1$ before standardization. By
\eqref{eq:ptilde}, $\widetilde{p}(s_t) \approx \mathbb{E}[R \mid s_t]$, so
$A^P$ is the empirical advantage that an idealized binary-reward critic
would compute: it is the change in the expected terminal reward attributed
to step $t$. This is the variant most directly compatible with the binary
verifier we use, but it suffers \emph{near saturation}: when $V_\phi(s_t)$
is large positive or negative, $\sigma'(V_\phi)$ collapses toward zero, so
two prefixes with very different log-ratios get essentially the same
gradient. In practice this means \emph{prob} undersignals on prefixes the
value model is highly confident about.

\item \emph{logprob}: The TD operates on the log-probability scale, mirroring
the form of token log-likelihoods used in the PPO ratio
$\rho_t(\theta)$. For a sequence model trained by maximum likelihood,
log-probability differences are the canonical update signal; using
$\log \widetilde{p}$ keeps the value-derived advantage on the same
information-theoretic scale as the underlying token logits. Numerically,
$\log \sigma(V) = -\log(1 + e^{-V})$, which behaves linearly for $V \ll 0$
(matching the value scale on the negative tail) and saturates softly to
zero for $V \gg 0$ (compressing the high-confidence positive tail). This is
a smooth interpolation between \emph{value} (linear in the negative regime)
and \emph{prob} (bounded above).
\end{itemize}

\paragraph{When each formulation is preferable.}
The three formulations are not interchangeable: each is the principled
choice under a different calibration regime of $V_\phi$.
\emph{prob} (\eqref{eq:advP}) is the natural statistical match when $V_\phi$
is well-calibrated as a Bradley--Terry score, and admits a clean
expected-terminal-reward interpretation $A^P(s_t) = r_t + \gamma\,
\mathbb{E}[R \mid s_{t+1}] - \mathbb{E}[R \mid s_t]$ via
\eqref{eq:ptilde}.
\emph{value} (\eqref{eq:advV}) is preferred when $V_\phi$ has wide log-ratio
dynamic range that the BCE objective does not directly bound, in which case
operating on $\sigma(V_\phi)$ would discard signal at the saturation tails;
the price is heavier-tailed advantages, controlled in our setup by
$\pm 3$ standardized clipping.
\emph{logprob} (\eqref{eq:advLP}) provides a smooth interpolation: it
preserves linear sensitivity in the negative-$V_\phi$ regime (where most
gradient signal lives early in training, when the policy is still mostly
off-distribution) and softly compresses the positive-$V_\phi$ regime
(where over-confident prefixes would otherwise dominate the batch). It also
matches the information-theoretic scale of the underlying token logits used
in the PPO ratio $\rho_t(\theta)$, which is a desirable invariance property
under reparameterization of $\pi_\phi$.
We treat these three as a contributed family of step-TD formulations and
report all three on MATH-500 pass@$16$ in Section~\ref{sec:results}.

\subsection{Step advantage normalization and clipping}
Let $\mathcal{A} = \{A(s_t)\}$ over all valid step boundaries in the outer
PPO step. We standardize over the entire batch
($\widetilde{A}_t = (A_t - \mu_{\mathcal{A}}) / (\sigma_{\mathcal{A}} +
\varepsilon)$) and clip the standardized advantages to $[-3, 3]$ to bound the
contribution of single-step outliers. The clip range is chosen empirically to
balance suppression of pathological prefixes against retaining genuine
step-level signal.

\section{Experimental Setup}
\label{sec:setup}

\subsection{Models and data}
We evaluate on two policy initializations and one shared dataset family:
\begin{itemize}
\item \textbf{Rho-1B}: \texttt{realtreetune/rho-1b-sft-MATH}, a 1.1B-parameter
SFT-on-MATH initialization used in the upstream \vineppo{} paper.
\item \textbf{DeepSeekMath-7B}: \texttt{realtreetune/deepseekmath-7b-sft-MATH-v2},
the 7B SFT-on-MATH initialization from the same upstream pipeline.
\end{itemize}
Training prompts are drawn from \texttt{DigitalLearningGmbH/MATH-lighteval}
(train split). Held-out evaluation uses \texttt{HuggingFaceH4/MATH-500} (test
split, 500 problems) plus the full MATH test set, College Math, and
OlympiadBench (math-only configuration). All settings use the prompt template
\texttt{[MATH\_TASK] Problem:\textbackslash n\{query\}\textbackslash n\textbackslash nSolution:}
and the LaTeX-aware step splitter from the upstream codebase.

\subsection{Methods}
We compare four methods:
\begin{description}
\item[\ppo{}] Terminal-reward, sequence-level advantages broadcast across
response tokens.
\item[\grpo{}] Per-prompt group-relative terminal-reward advantages with
$G=8$.
\item[\vineppo{}] $K=9$ Monte Carlo continuations from each non-terminal step
boundary; step-TD advantages.
\item[\caspo{}] \ipvrm{} prefix value model with online updates; step-TD
advantages over $V_\phi$. We report all three step-TD formulations from
Section~\ref{sec:transforms} (\caspo{}-\emph{value}, \caspo{}-\emph{prob},
\caspo{}-\emph{logprob}) as separate runs, and a fourth variant
\caspo{}-frozen-RM that uses the \emph{value} formulation but disables
online updates of $V_\phi$ during policy training.
\end{description}

\subsection{Hyperparameters}
All four methods share the same PPO infrastructure. Table~\ref{tab:hp}
summarizes the key knobs.

\begin{table}[t]
\centering
\caption{Shared training hyperparameters (Rho-1B and 7B). All methods use the
same PPO clipped objective, KL coefficient, group size, response budget, and
global PPO minibatch.}
\label{tab:hp}
\begin{tabular}{lll}
\toprule
Field & Rho-1B & DeepSeekMath-7B \\
\midrule
Group size $G$ & 8 & 8 \\
Prompts per outer step & 64 & 8 \\
Responses per outer step & 512 & 64 \\
Global PPO minibatch & 64 & 64 \\
PPO epochs per rollout & 2 & 2 \\
Policy LR & $1\times 10^{-6}$ & $2\times 10^{-6}$ \\
Online value LR (\caspo{}) & $1\times 10^{-6}$ & $1\times 10^{-6}$ \\
Value $\beta$, $m$ & $10$, $5$ & $10$, $5$ \\
PPO clip $\epsilon$ & $0.2$ & $0.2$ \\
KL coefficient & $10^{-4}$ ($k_3$) & $10^{-4}$ ($k_3$) \\
Advantage clip & $3.0$ & $3.0$ \\
Advantage standardization & batch & batch \\
Rollout sampling & temp $0.6$, top-$p$ $0.9$ & temp $0.6$, top-$p$ $0.9$ \\
Max prompt / response & 1024 / 1024 & 1024 / 1024 \\
Outer steps & 1000 & 1500 \\
Checkpoints & step\_250/500/750/final & step\_250/500/750/1000/1250/final \\
\bottomrule
\end{tabular}
\end{table}

The 7B run uses 1.5$\times$ the outer-step budget of the Rho-1B run because
the 7B base policy needs more updates per equivalent training signal. The
larger Rho-1B prompt batch reflects its lower per-step compute cost; both
configurations land at a 64-response global PPO minibatch.

\subsection{Evaluation protocol}
Final evaluation reports pass@1 and pass@$k$ ($k=16$) on MATH-500, full MATH
test, College Math, and OlympiadBench. Sampling at evaluation uses
temperature $0.35$ and top-$p$ $0.9$, matching the upstream protocol.
Generation is performed with vLLM at \texttt{kv\_cache\_dtype="fp8"} and
\texttt{gpu\_memory\_utilization=0.92} for maximum eval throughput on a
dedicated GPU. Sample evaluation at intermediate checkpoints uses
\texttt{EVAL\_LIMIT=100}, \texttt{EVAL\_K=8} for low-overhead training-curve
visibility; the full benchmark suite is run at the final checkpoint.

\subsection{Infrastructure}
Both scales share a single PyTorch trainer plus rank-local vLLM rollout. The
1B configuration uses a single H100 per method; the 7B configuration uses
FSDP=4 (full\_shard, use\_orig\_params) for \ppo{}/\grpo{}/\caspo{} and FSDP=8
for \vineppo{}, with vLLM colocated rank-local in both layouts. Trainer-to-vLLM
weight synchronization uses CUDA IPC: at 1B with a single-rank trainer, the
IPC handle exchange is direct; at 7B under FSDP, the trainer enters
\texttt{FSDP.summon\_full\_params(writeback=False)} collectively before
issuing per-tensor IPC pushes through each rank's vLLM
\texttt{EngineCore.collective\_rpc("update\_weights", \dots)}.

For numerical stability under the FSDP layout, the trainer co-trains policy
and (for \caspo{}) value model in bf16 with fp32 master weights; the
reference policy is shared between the KL term and \caspo{}'s value-model ref
to avoid a duplicate forward. Rollout, old-logprob rescore, ref-logprob
precompute, and policy/value forward use FlashAttention-3 on Hopper.

\section{Results}
\label{sec:results}

This section will report wall-clock and accuracy after the full training and
evaluation runs complete. The intended structure is:

\begin{itemize}
\item Per-method MATH-500 pass@1 and pass@$16$ at \texttt{final}, with
seed-level error bars over a small number of seeds (target: 3 seeds per
method per scale).
\item Per-method full-suite (MATH, College Math, OlympiadBench) pass@$16$ at
\texttt{final}.
\item Training curves on MATH-500 sample evaluation at \texttt{step\_250},
\texttt{step\_500}, \texttt{step\_750}, and \texttt{final}, summarizing
sample efficiency.
\item Per-formulation comparison of the three step-TD advantage variants
(\caspo{}-\emph{value} vs.\ \caspo{}-\emph{prob} vs.\ \caspo{}-\emph{logprob})
on MATH-500 pass@$16$, mapping the empirical winner to the calibration
regime predicted in Section~\ref{sec:transforms}.
\item \caspo{}-frozen-RM vs.\ \caspo{} (default, online value updates on)
to quantify the contribution of the online update path.
\item Wall-clock summary: per-method step time, per-method 1500-step ETA,
and the implied GPU-hour cost of the four-method comparison.
\end{itemize}

We report all numbers under the same training contract: identical SFT
initialization, identical prompt set, identical verifier, identical
group/response budget, identical PPO clipped objective, identical KL
coefficient, identical advantage standardization, identical max sequence
length, and identical evaluation sampling parameters. The only intended
difference between rows is the source of the per-step value $V(s_t)$ in
\eqref{eq:adv}: zero (\ppo{}/\grpo{} terminal broadcast), Monte Carlo
continuations (\vineppo{}), or learned prefix value $V_\phi$ (\caspo{}).

\section{Analysis}
\label{sec:analysis}

This section will examine, given the empirical results, what determines when
$V_\phi$ matches \vineppo{}'s prefix-rollout signal and when it does not. The
intended structure is:

\begin{itemize}
\item \emph{Step-credit calibration}: do step advantages from $V_\phi$ rank
correct vs.\ incorrect step continuations consistently with on-policy MC?
\item \emph{Online value drift}: does the online update preserve the
phase-1 calibration, or does it collapse to constant predictions? We monitor
$\bar V^+$ and $\bar V^-$ (mean prefix value at correct vs.\ incorrect step
continuations) over outer steps to detect drift.
\item \emph{Wall-clock per accuracy point}: the practical question \caspo{}
is asked to answer is whether learned prefix values can match \vineppo{}'s
accuracy at a fraction of the wall-clock cost. We report this as
accuracy-versus-GPU-hours for each method.
\item \emph{Step-count effects on \vineppo{}}: \vineppo{}'s wall-clock
scales as $K \cdot \text{steps/response}$, so the marginal cost of step-level
credit grows with response length. We quantify this by binning the test set
into long vs.\ short reasoning regimes and reporting the per-bin
accuracy/wall-clock trade-off.
\end{itemize}

\section{Limitations}
\label{sec:limitations}

\paragraph{Reward-model drift.}
\caspo{}'s online value updates are gated by ADB and DLW corrections, but no
correction perfectly defends against pathological feedback loops (e.g.\ a
value head that learns to predict the policy's surface form rather than its
correctness). The \caspo{}-frozen-RM variant in our contributed family
isolates this risk by holding $V_\phi$ fixed during policy training, at the
cost of distribution-shift bias.

\paragraph{Segmenter quality.}
The LaTeX-aware step splitter is a structural heuristic. Manual spot checks
on Rho-1B MATH generations show coherent boundaries for prose, equations,
factorization, and aligned derivations, but generations that continue past a
final boxed answer have those trailing fragments treated as additional steps;
a post-answer truncation pass is left as future work.

\paragraph{Apples-to-apples scope.}
Our comparison fixes the SFT initialization, prompt set, verifier, sampling
parameters, KL constraint, and advantage standardization. It does not
exhaustively explore method-specific tuning (e.g.\ different $K$ for
\vineppo{}, different $\beta/m$ for \caspo{}, larger $G$ for \grpo{}); we
report the upstream-paper values where available and otherwise the values
that minimize wall-clock without regressing pass@1 on a small validation set.

\paragraph{Scale.}
Results at 1B and 7B may not transfer to substantially larger scales. The
7B FSDP+vLLM-IPC stack is novel to this work; whether the same configuration
extrapolates to 70B is an open systems question.

\section{Conclusion}
\label{sec:conclusion}

We described \caspo{}, a step-level credit assignment method for reasoning
RL that replaces \vineppo{}'s on-policy Monte Carlo prefix rollouts with an
\ipvrm{}-style learned prefix value model and an optional online value
update. \caspo{} fits cleanly into the same PPO clipped objective, the same
KL constraint, and the same step-TD advantage construction used by
\vineppo{}, isolating the credit-signal change from confounders. By unifying
\ppo{}, \grpo{}, \vineppo{}, and \caspo{} under one trainer, one verifier,
and one rollout engine, we make it possible to compare these four methods
under identical conditions on Rho-1B and DeepSeekMath-7B. Empirical results
in Section~\ref{sec:results} will quantify whether the learned value model
recovers \vineppo{}'s reasoning-credit signal at a fraction of its
generation cost, or whether step-level Monte Carlo remains uniquely
advantageous.

\bibliographystyle{plainnat}
\bibliography{references}

\end{document}
